% Appendix C — Robustness Specifications
% Drafted per skeleton v5.1 §V flag (BelowIG handling + OPSW log-DV future-work).

\section{Robustness Specifications}
\label{app:robustness}

This appendix collects auxiliary specifications and design-choice disclosures supporting the main \S\ref{sec:main} hypotheses HC and HFC.

\subsection{Three-Tier Constraint Moderator: BelowIG Handling}
\label{app:robustness:belowig}

The HFC specification of \S\ref{sec:main:hfc} estimates the financial-constraint moderator using a three-tier rating partition (IG, BelowIG, Unrated) but reports the IG-versus-Unrated contrast as the main-display estimand. The BelowIG category enters the model as a level shift relative to the IG reference category, with no BelowIG-specific interaction term constructed in the H1.2.ceo2 specification. This restriction reflects two design considerations. First, the FP framework on which HFC is built is binary (rated vs.\ unrated); the three-tier extension is our refinement, and within the rated half of the population the BelowIG cell is small and economically heterogeneous. Second, the main hypothesis localizes the precautionary cash response to the credit-constrained Unrated cell as predicted by FP, so the displayed contrast is between the credit-constrained Unrated and the rated pool (IG $\cup$ BelowIG combined as the reference). The full BelowIG level-coefficient row is reported in the body Table~\ref{tab:h1_2_ceo2}; readers wanting a fully decomposed three-tier interaction specification can re-estimate H1.2 with an explicit BelowIG~$\times$~UncAnsCEO\_c term added to the model.

\subsection{Alternative Cash Dependent-Variable Specification (OPSW Log-Form)}
\label{app:robustness:opsw}

The body adopts the linear cash-to-assets dependent variable (\texttt{cheq/atq}) of \citeA{bates2009} as the standard precautionary-cash form. \citeA{opler1999} also use a log-form dependent variable, defined as $\log(\text{cheq}/\text{atq})$ (or, in some specifications, $\log(\text{cheq}/(\text{atq}-\text{cheq}))$ to avoid log-of-zero issues at the lower tail). The log form down-weights the upper tail of the cash distribution relative to the linear form and tests whether the precautionary association is driven by high-cash firms.

The present draft does not estimate the OPSW log-DV robustness; it is flagged in \S\ref{sec:conclusion} as a planned extension. The expectation, given the magnitude of the linear-form coefficients in \S\ref{sec:main:hc} (UncAnsCEO sig 12/12 at $p<0.10$), is that the qualitative result survives under the log specification, with attenuated point-estimate magnitudes due to the tail-compression. Future work will report the log-DV coefficients directly to confirm this expectation.

\subsection{Fixed-Effects Ladder Robustness}
\label{app:robustness:fe}

The four-step fixed-effects ladder of \S\ref{sec:framework:design} (Industry$+$Year, Firm$+$Year, Industry$+$Year-Quarter, Firm$+$Year-Quarter) is the within-regressor robustness frame applied throughout the body. Every column of every body table reports a different point on this ladder, so the FE-ladder robustness check is integrated rather than a separate appendix specification. The convergence pattern across the ladder---attenuation from industry to firm fixed effects, additional attenuation from year to year-quarter fixed effects---is the standard within-firm panel-OLS attenuation pattern (within-firm variation is a strict subset of cross-firm variation, so coefficients on a regressor with substantial across-firm dispersion shrink as fixed effects absorb that dispersion). The signs and significance of UncAnsCEO are robust to all four ladder steps in HC and HFC; this robustness is documented in the body and not separately tabulated here.
